\begin{pastexam}[label={px:ntu-109-ee-ds-04}]{臺大電機所 109 學年度（B）第 4 題}
（單選；每題 4 分，答錯不倒扣）

Please choose the best possible answer in each of the following questions. In questions that ask
for time complexities, please choose the tightest bound.

The worst case time complexity of one union operation upon two binary heaps of total size n is

(A)~O(1) \qquad (B)~O(log*n) \qquad (C)~O(log n) \qquad (D)~O(n) \qquad (E)~O(n log n)
\end{pastexam}

\begin{pastsoln}
答案是 (D) O(n)。

二元堆積存在陣列裡，形狀固定是完整二元樹，兩個堆積沒辦法像左傾樹或二項堆積那樣只改幾個指標就
接起來（資結書第~8.5~節）。 合併的做法是把兩個陣列接成一個長度 $n$ 的陣列，再用由下而上的建堆積
重建一次，$O(n)$（分析見歷屆試題~\ref{px:ntu-109-ee-ds-02}）。

下界：結果要放在同一個陣列裡，至少得把其中一個堆積的元素全部搬進另一個陣列。 兩個堆積各有
$n/2$ 個元素時，光搬動就是 $\Omega(n)$。 所以最緊的界是 $O(n)$。

另一種做法是把小的那個堆積的元素一個一個插入大的，最壞是 $\Theta(n \log n)$，比重建慢。
(A) O(1) 是費氏堆積（平攤）的合併時間，(C) O(log n) 是左傾樹和二項堆積的合併時間，都要靠指標結構，
二元堆積做不到；(B) 比 (C) 還小，同樣做不到；(E) 是上界但不緊。
\end{pastsoln}
